\section*{Chapitre \ref{ch:iv2:markets-and-microstructure} --- Marchés et microstructure}
\markboth{\protect\omTocMark{}{Réponses~: Marchés et microstructure}}{\protect\omTocMark{}{Réponses~: Marchés et microstructure}}

\begin{solution}{iq:iv2:markets-and-microstructure:1}
Le prix moyen obtenu est d'environ 100.0267, soit environ 2.67 cents au-dessus du prix milieu de 100.00. Le meilleur cours vendeur devient 100.03, pour 200 actions, et le prix milieu devient 100.005.

On prend les 100 à 100.02 et 200 des 400 à 100.03, donc le total vaut $10\,002 + 20\,006 = 30\,008$, soit 100.0267 par action. Le prix milieu avant était $(99.98 + 100.02)/2$. Ensuite, il vaut $(99.98 + 100.03)/2$.

\iqwrong{Utiliser le meilleur cours vendeur pour les 300 actions. Seules les 100 affichées sont disponibles à 100.02.}
\end{solution}

\begin{solution}{iq:iv2:markets-and-microstructure:2}
Le prix milieu est 100.01. Le microprix\index{microprix} vaut 100.018. On s'attend à ce que le prochain mouvement soit à la hausse, car peu de taille se trouve au cours vendeur.

Le microprix est $(100.00 \times 100 + 100.02 \times 900)/1000$. Le prix dont la file est la plus courte est plus proche d'être traversé, si bien que l'estimation de la fair value penche de son côté. On notera le croisement des pondérations~: la grande taille à l'achat tire l'estimation vers le cours vendeur.

\iqwrong{Pondérer chaque prix par sa propre taille, ce qui donne 100.002 et penche dans le mauvais sens.}
\iqfollow{Une grosse taille à l'achat est-elle toujours une bonne nouvelle~?}{Non. La taille affichée peut être annulée, et sur un marché peu liquide elle peut être là pour induire en erreur. Il faut donc la traiter comme une statistique à tester, et non comme un fait.}
\end{solution}

\begin{solution}{iq:iv2:markets-and-microstructure:3}
Avec la priorité prix-temps, $A$ obtient 100, $B$ obtient 150, et $C$ n'obtient rien. Avec le pro-rata, ils obtiennent environ 71.4, 142.9 et 35.7. La priorité temporelle récompense la rapidité, et le pro-rata récompense la taille.

Le total des ordres est 350. Les parts au pro-rata valent $250 \times 100/350$, $250 \times 200/350$ et $250 \times 50/350$. En pratique, elles sont arrondies à des lots entiers, avec une règle pour le reliquat.

Avec le pro-rata, afficher plus que ce que l'on souhaite permet d'obtenir davantage d'une petite transaction, si bien que les intervenants surdimensionnent leurs cotations. Avec la priorité temporelle, la valeur se trouve en tête de file, et les acteurs se disputent la place pour y parvenir.

\iqwrong{Croire que les règles de priorité n'intéressent que les courtiers. Elles déterminent de quoi se compose l'avantage d'un teneur de marché.}
\end{solution}

\begin{solution}{iq:iv2:markets-and-microstructure:4}
La fourchette effective\index{fourchette effective} vaut 6 cents, la fourchette réalisée\index{fourchette réalisée} 2 cents, et l'impact sur le prix 4 cents. Le teneur de marché a conservé la fourchette réalisée.

La fourchette effective est égale au double de l'écart au prix milieu au moment de la transaction~: $2 \times (100.03 - 100.00)$. La fourchette réalisée est égale au double de l'écart au prix milieu ultérieur~: $2 \times (100.03 - 100.02)$. L'impact sur le prix est égal au double du déplacement du prix milieu~: $2 \times 0.02$. Les deux parties s'additionnent~: $2 + 4 = 6$.

On a payé 3 cents. Le teneur de marché a vendu à 100.03 ce qui vaut 100.02 cinq minutes plus tard. Il ne lui reste donc qu'un cent sur ses 3, et les 2 autres sont allés à la sélection adverse.

\iqwrong{Prendre la fourchette effective pour le profit du teneur de marché. Une grande partie peut en être perdue à cause du mouvement du prix.}
\end{solution}

\begin{solution}{iq:iv2:markets-and-microstructure:5}
On économise 1.5 cents par action. Un teneur de marché dégage 1.4 cents net sur un aller-retour.

Traverser la fourchette coûte 1 cent et les frais de 0.3 cent, soit 1.3 cents de plus que le cours acheteur. Une exécution passive au cours acheteur rapporte la rémunération de 0.2 cent. La différence est de 1.5 cents. Un teneur de marché qui achète au cours acheteur et vend au cours vendeur gagne 1 cent et deux rémunérations~: $1 + 2 \times 0.2 = 1.4$.

Il s'agit d'une plateforme maker-taker\index{maker-taker}. Les frais sont en partie inclus dans le cours affiché, si bien que la fourchette visible n'est pas ce que paient réellement les deux parties. Une fois les frais déduits, la bourse conserve 0.1 cent.

\iqwrong{Supposer qu'on obtient toujours le gain. Un ordre passif s'exécute moins souvent, et il s'exécute justement quand le prix est sur le point de traverser l'ordre.}
\end{solution}

\begin{solution}{iq:iv2:markets-and-microstructure:6}
À 5~\$, un cent représente 20 points de base, et à 500~\$, 0.2 point de base. L'action à 5~\$ a davantage de chances de rester bloquée sur un seul tick avec une longue file. Là, la position dans la file compte davantage que le prix, puisque personne ne peut améliorer le prix.

Si le spread naturel est inférieur au tick, disons 5 points de base, l'action à 5~\$ ne peut pas coter ce niveau. Tout le monde se place donc au même niveau, et la file est longue. Le spread de l'action à 500~\$ peut varier avec peu de contraintes, si bien que son carnet est peu fourni à chacun des nombreux niveaux de prix, et la concurrence par les prix joue.

\iqwrong{Croire que les actions à bas prix sont peu coûteuses à négocier. En pourcentage, un spread contraint par le tick peut être bien supérieur au spread juste.}
\end{solution}

\begin{solution}{iq:iv2:markets-and-microstructure:7}
Le spread vaut 4 cents, d'après $s = 2\sqrt{-\mathrm{Cov}} = 2\sqrt{0.0004}$. L'autocovariance est négative parce qu'un achat au cours vendeur tend à être suivi d'une transaction au cours acheteur, ce qui fait revenir le prix.

C'est l'estimateur de Roll\index{modèle de Roll}. Le prix vaut $p_t = m_t + \tfrac{s}{2}q_t$, avec $q_t = \pm1$ indépendants. La variation vaut $\Delta m_t + \tfrac{s}{2}(q_t - q_{t-1})$. Seul le terme en $q_{t-1}$ est commun à deux variations successives, si bien que la covariance vaut $-(s/2)^2\,\mathrm{Var}(q) = -s^2/4$.

Sans rebond, $s = 0$, et la covariance est nulle, comme pour une marche aléatoire.

\iqwrong{L'utiliser lorsque l'autocovariance est positive. La formule n'a alors pas de réponse réelle, ce qui montre que les hypothèses ne tiennent pas~: les transactions se regroupent, ou le prix suit une tendance.}
\end{solution}

\begin{solution}{iq:iv2:markets-and-microstructure:8}
Le prix moyen pondéré par les volumes\index{VWAP} vaut 10.11. On le bat de 3 cents, soit environ 30 points de base. La moyenne simple des trois prix vaut 10.10.

La somme des prix multipliés par les volumes vaut $1\,000 + 2\,040 + 7\,070 = 10\,110$, pour 1\,000 actions. Le prix obtenu, 10.08, est 3 cents en dessous, et $0.03/10.11 \approx 0.30~\%$.

L'indice de référence est donc pondéré par les niveaux où le volume a réellement été échangé. La moyenne simple donne aux petits volumes à 10.00 le même poids qu'au gros volume à 10.10.

\iqwrong{Croire que battre l'indice de référence prouve une bonne exécution. Un acheteur qui a commencé tôt peut simplement avoir négocié avant que son propre impact ne se fasse sentir.}
\end{solution}

\begin{solution}{iq:iv2:markets-and-microstructure:9}
Les jours de couverture\index{jours de couverture} sont de 15. À 20~\% du volume, il faudrait 75 jours. Une action qui bondit peut contraindre les vendeurs à découvert à racheter à n'importe quel prix, et une longue file de sortie accroît leur urgence.

Le ratio vaut $30/2 = 15$. Acheter 20~\% de 2 millions représente 0.4 million par jour, et $30/0.4 = 75$.

Ce chiffre mesure l'ampleur d'une demande forcée possible face à la liquidité normale. Les vendeurs à découvert dont les stops sont déclenchés, ou dont le prêteur rappelle le titre, doivent racheter à n'importe quel prix, ce qui fait monter le prix. Cela déclenche alors les stops d'autres vendeurs. C'est un short squeeze\index{short squeeze}.

\iqwrong{Dire qu'un encours vendeur à découvert élevé signifie que l'action doit monter. Les vendeurs peuvent avoir raison. Il indique l'ampleur possible du mouvement s'ils ont tort.}
\end{solution}

\begin{solution}{iq:iv2:markets-and-microstructure:10}
La fair value est de 25.00~\$. Le certificat de dépôt est surcoté de 80 points de base~: il faut donc le vendre, acheter les actions et couvrir la devise, ce qui laisse environ 50 points de base, soit environ 12.5 cents par certificat.

Deux actions coûtent $\pounds 20$, et à 1.25, cela fait 25~\$. On a alors $0.20/25 = 0.80~\%$, et $80 - 30 = 50$ points de base. La conversion entre certificat et actions est ce qui lie les deux prix, et elle coûte des frais.

\iqwrong{Oublier la devise. Si l'on achète les actions en livres et que l'on vend le certificat en dollars, une position subsiste tant qu'on ne la clôture pas ou qu'on ne la couvre pas.}
\end{solution}

\begin{solution}{iq:iv2:markets-and-microstructure:11}
Pas sous le régime de protection des ordres~: il faut d'abord acheter les 100 à 100.01 sur $A$, puis les 200 à 100.02 sur $B$. Un trade-through\index{trade-through} est une transaction à un prix moins bon qu'une cotation protégée sur une autre plateforme. Un ordre de balayage intermarchés (ISO) permet d'envoyer vers $B$ immédiatement, à condition de promettre d'envoyer vers $A$ au même moment.

La Rule 611 de la Regulation NMS\index{Regulation NMS} oblige les plateformes à en assurer le respect. L'ordre de balayage intermarchés est l'outil de routage~: l'ordre est marqué pour indiquer que l'on a également consommé les meilleurs prix, et la plateforme n'a pas à l'acheminer ailleurs.

\begin{dated}{2026-11}{Protection des ordres en cours de révision}\label{dat:iv2:markets-and-microstructure:rule611}
La Rule 611 a été adoptée en 2005. Le 11 juin 2026, la SEC a proposé de l'abroger, ainsi que la Rule 610(e). Il faut vérifier où en est ce dossier avant d'énoncer la règle comme étant en vigueur.
\end{dated}

\iqwrong{Croire que la meilleure plateforme est simplement celle où l'on obtient la meilleure exécution aujourd'hui. L'obligation de routage fixe où l'on peut négocier en premier.}
\end{solution}

\begin{solution}{iq:iv2:markets-and-microstructure:12}
Le prix est 100, et 300 actions s'échangent. À 100, tous les acheteurs à 100 ou plus sont exécutés, de même que tous les vendeurs à 100 ou moins. Les acheteurs à 99 et les vendeurs à 101 ne sont pas exécutés.

On compare, à chaque prix, la demande et l'offre cumulées~:
\begin{center}\small
\begin{tabular}{@{}rrrr@{}}
Prix & Acheteurs prêts & Vendeurs prêts & Volume \\ \hline
99 & 600 & 150 & 150 \\
100 & 300 & 300 & 300 \\
101 & 100 & 500 & 100
\end{tabular}
\end{center}
À 100, les deux côtés valent exactement 300. C'est une enchère\index{enchère}~: tout le monde échange à un prix unique.

\iqwrong{Apparier les ordres un à un depuis le haut, comme sur un marché continu. L'intérêt de l'enchère est que tous obtiennent le même prix.}
\end{solution}

\begin{solution}{iq:iv2:markets-and-microstructure:13}
On peut exécuter environ 222\,222 actions, de sorte que l'ordre est terminé aujourd'hui. Cela représente 11.1~\% du volume des autres.

Si l'on représente 10~\% du total, alors $x = 0.1(2\,000\,000 + x)$, donc $x = 200\,000/0.9$. Ainsi, les 200\,000 actions peuvent être exécutées lorsque les autres ont échangé 1.8~million d'actions.

Un ordre au pourcentage du volume (POV) s'exécute davantage lorsque le volume est élevé, et le volume est souvent élevé quand le cours s'emballe. Il a donc tendance à acheter après les hausses, et ses propres transactions alimentent le volume qu'il suit ensuite.

\iqwrong{Prendre 10~\% des 2~millions des autres, soit 200\,000. Le taux doit inclure les propres transactions de l'ordre, sinon celui-ci s'exécute au mauvais rythme.}
\end{solution}

\begin{solution}{iq:iv2:markets-and-microstructure:14}
Le courtier doit avoir des motifs raisonnables de croire que les actions peuvent être empruntées pour la livraison~: c'est l'exigence de locate\index{locate}. Un défaut de livraison survient lorsque le vendeur ne livre pas à la date de règlement-livraison. Une vente à découvert sans emprunt n'est pas nécessairement illégale. Les teneurs de marché de bonne foi bénéficient d'une dérogation à l'exigence de locate.

Dans le cadre de la Regulation SHO\index{Regulation SHO}, la Rule 203(b)(1) fixe l'exigence de locate, et la Rule 204 fixe le moment où un défaut doit être clôturé, au début des négociations le lendemain de la date de règlement-livraison, pour une vente à découvert. Un teneur de marché doit pouvoir vendre avant d'avoir trouvé un emprunt, afin de coter. La dérogation ne couvre pas les ventes spéculatives.

\iqwrong{Assimiler un défaut de livraison à une fraude. Les défauts surviennent pour des raisons opérationnelles ordinaires. Les règles fonctionnent en imposant une clôture forcée.}
\end{solution}

\begin{solution}{iq:iv2:markets-and-microstructure:15}
1~\% du volume d'une journée coûte 20~points de base, et 4~\% coûte 40~points de base. Quadrupler l'ordre multiplie par 8 le coût en dollars.

$2~\% \times \sqrt{0.01} = 0.2~\%$, et $2~\% \times \sqrt{0.04} = 0.4~\%$. Le coût par action double, et il y a quatre fois plus d'actions. Le coût total croît donc comme $Q^{3/2}$.

C'est la loi en racine carrée\index{loi d'impact en racine carrée} sous une forme simplifiée~: une régularité empirique de l'impact moyen, avec une constante proche de 1 dans de nombreuses études. Elle s'applique aux ordres bien inférieurs au volume d'une journée entière, et il s'agit d'une moyenne. Un ordre donné peut faire nettement mieux ou nettement pire.

\iqwrong{Supposer un impact linéaire, de sorte que quatre fois la taille signifie quatre fois le coût par action. Cela surestime les grands ordres, mais le coût total reste plus que proportionnel.}
\iqwhy{Elle fixe la capacité d'une stratégie~: le coût par dollar croît avec la racine carrée de la taille, de sorte qu'à partir d'une certaine taille il absorbe l'avantage (edge).}
\end{solution}

\begin{solution}{iq:iv2:markets-and-microstructure:16}
L'implementation shortfall\index{implementation shortfall} vaut 1\,480~\$, soit 29.6 points de base du portefeuille fictif de 500\,000~\$. Il se décompose en 800~\$ de coût d'exécution, 80~\$ de commission et 600~\$ de coût d'opportunité.

L'exécution vaut $8\,000 \times \$0.10$. La commission vaut $8\,000 \times \$0.01$. Les 2\,000 actions non achetées auraient rapporté $\$0.30$ chacune. Une vérification par le portefeuille fictif~: le rendement fictif vaut $10\,000 \times 0.30 = 3\,000$. Le rendement réel vaut $8\,000 \times 0.20 - 80 = 1\,520$. L'écart vaut 1\,480.

\iqwrong{Oublier les actions non achetées. Un trader qui n'achète rien n'a aucun coût d'exécution, mais il a manqué tout le mouvement.}
\end{solution}

\begin{solution}{iq:iv2:markets-and-microstructure:17}
On est exécuté au 21\up{e} événement. En moyenne, cela prend 14 secondes, avec un écart type d'environ 3.06 secondes. La probabilité d'être exécuté en moins de 20 secondes est d'environ 96~\%.

L'attente est la somme de 21 intervalles exponentiels, soit une loi gamma, de moyenne $21/1.5$ et d'écart type $\sqrt{21}/1.5$. Être exécuté d'ici 20 secondes signifie qu'il y a au moins 21 événements en 20 secondes, et ce nombre suit une loi de Poisson de moyenne 30.

\iqwrong{Utiliser 20 événements. Il faut que les 20 ordres devant soient retirés, puis qu'un ordre supplémentaire soit exécuté.}
\iqfollow{Qu'est-ce que le modèle laisse de côté~?}{Les annulations surviennent n'importe où dans la file, et les taux évoluent avec le marché. Le plus important est que l'exécution arrive lorsque la file s'effondre, souvent au moment où le cours traverse le prix de l'ordre.}
\end{solution}

\begin{solution}{iq:iv2:markets-and-microstructure:18}
En tête de file, on espère gagner 0.3 cent par action, et en queue de file on perd 0.2 cent. La valeur de la position dans la file est donc de 0.5 cent par action dans ce modèle.

On gagne une demi-fourchette par rapport au prix milieu, soit 0.5 cent. La baisse est d'un tick, soit 1 cent, de sorte que son coût espéré vaut 0.2 cent en tête de file et 0.7 en queue. Ainsi, $0.5 - 0.2 = 0.3$ et $0.5 - 0.7 = -0.2$.

C'est pourquoi une position dans la file a une valeur, et pourquoi les firmes investissent dans la vitesse pour atteindre la tête de file. La même cotation, au même prix, est une bonne transaction à un endroit de la file et une mauvaise à un autre.

\iqwrong{Traiter toutes les exécutions à un prix comme identiques. La différence entre elles tient à celui avec qui l'on négocie.}
\end{solution}

\begin{solution}{iq:iv2:markets-and-microstructure:19}
Le cours vendeur (ask) est d'environ 108.33 et le cours acheteur (bid) d'environ 104.09. Ils ne sont pas symétriques autour de la moyenne de 106, car la probabilité a priori n'est pas un demi. Après un achat et une vente, la probabilité de 120 revient à 30~\%.

Un achat est plus probable si la valeur est élevée~: $P(\text{achat} \mid 120) = 0.25 + 0.375 = 0.625$, et $P(\text{achat} \mid 100) = 0.375$. Par la formule de Bayes\index{modèle de Glosten--Milgrom}, $P(120 \mid \text{achat}) = 0.1875/(0.1875 + 0.2625) \approx 0.4167$, et le cours vendeur vaut $100 + 20 \times 0.4167 \approx 108.33$. De même, $P(120 \mid \text{vente}) = 0.1125/(0.1125 + 0.4375) \approx 0.2045$, et le cours acheteur vaut environ 104.09.

Le cours vendeur est supérieur de 2.33 à la moyenne, et le cours acheteur inférieur de 1.91. Un achat est moins probable lorsque la valeur est plus probablement basse, de sorte qu'il porte davantage d'information. Après un achat et une vente, le rapport de vraisemblance vaut $(0.625/0.375)(0.375/0.625) = 1$.

\iqwrong{Coter de façon symétrique autour de 106. Chaque cotation doit être l'espérance sachant qu'elle est touchée, et les deux montants d'information diffèrent.}
\iqfollow{Et si la part d'informés monte à 100~\%~?}{Alors un achat signifie 120 et une vente signifie 100. Le cours acheteur et le cours vendeur valent 100 et 120, et aucun trader ordinaire ne peut négocier à l'intérieur.}
\end{solution}

\begin{solution}{iq:iv2:markets-and-microstructure:20}
Le markout\index{markout} à une seconde vaut $+0.7$ cent par action, et à une minute $-0.3$ cent. Sur un million d'actions par jour, cela représente 7\,000~\$ et $-\$3\,000$. L'horizon pertinent est celui auquel on peut réellement dénouer sa position, et à une minute on a de fortes chances d'être en perte.

Le profit valorisé au prix milieu ultérieur $m$ vaut $m - 10.00$, plus le rebate. À une seconde, $0.5 + 0.2 = 0.7$ cent. À une minute, $-0.5 + 0.2 = -0.3$.

La chute, de 1 cent à 0.5 puis à $-0.5$, est la sélection adverse qui se manifeste. Il faut donc vérifier la durée réelle de détention. Il faut ensuite élargir la cotation face au flux qui coûte, ou sortir plus vite, ou cesser de coter à cet endroit de la file.

\iqwrong{Présenter la fourchette captée à l'exécution, soit 1.2 cent, comme profit du desk. C'est le chiffre qui a le meilleur aspect et qui signifie le moins.}
\end{solution}

\begin{solution}{iq:iv2:markets-and-microstructure:21}
Cela coûte 75~\$~: 25~\$ de financement net et 50~\$ de dividende. Le titre doit baisser de 75 cents pour atteindre 49.25~\$. La perte n'a pas de limite supérieure.

Le produit de la vente vaut 5\,000~\$. Sur un trimestre, on perçoit $4~\% \times 5\,000/4 = 50$ et on verse $6~\% \times 5\,000/4 = 75$, soit un net de 25. Il faut verser au prêteur le dividende, $100 \times 0.50 = 50$. Chaque dollar de baisse du titre rapporte 100~\$, de sorte que 75~\$ correspond à 75 cents par action.

À titre de vérification, si les frais égalent les intérêts et qu'il n'y a pas de dividende, le coût est nul.

\iqwrong{Oublier le dividende, ou penser qu'une vente à découvert ne rapporte rien sur son produit. Le taux net est ce qui compte, et pour un titre difficile à emprunter, il est fortement négatif.}
\end{solution}

\begin{solution}{iq:iv2:markets-and-microstructure:22}
Une baisse de 10~\% ou plus par rapport à la clôture de la veille, ici jusqu'à 45~\$ ou moins, déclenche le test. Pour le reste de la journée et toute la journée du lendemain, on ne peut vendre à découvert qu'au-dessus du meilleur cours acheteur national, soit à 44.89~\$ ou plus à présent. Si l'on détient les actions, on peut toujours frapper le cours acheteur.

La Rule 201 de la Regulation SHO\index{Regulation SHO} est la règle alternative de l'uptick. Elle n'interdit pas la vente à découvert, mais elle empêche les ventes à découvert de peser sur le cours acheteur~: les ordres à découvert doivent reposer dans le carnet d'ordres. Un système d'exécution doit donc connaître son état pour chaque titre, et savoir si chaque vente est marquée long ou short.

\iqwrong{Dire que la vente à découvert est interdite tant que le test est en vigueur. Seul le prix de la vente est limité.}
\end{solution}

\begin{solution}{iq:iv2:markets-and-microstructure:23}
L'indice de référence vaut 100.285. La moyenne pondérée par le temps vaut 100.425, soit environ 14 points de base au-dessus. Un calendrier d'exécution qui suit le volume aurait reproduit exactement l'indice de référence.

La moyenne pondérée par les volumes vaut $0.40 \times 100 + 0.15 \times 100.5 + 0.15 \times 101 + 0.30 \times 100.2 = 100.285$. La moyenne simple vaut 100.425. Son écart de 14~cents résulte de la sous-pondération de la première période, bon marché.

À la limite, si les quatre cours étaient égaux, les deux moyennes seraient égales et l'écart nul. Le choix du calendrier ne compte que dans la mesure où cours et volumes varient ensemble.

\iqwrong{Croire que cela prouve qu'un calendrier suivant le volume est moins cher. Cela établit seulement qu'il suit l'indice de référence. Un autre jour, la première période est celle où l'on paie cher.}
\end{solution}

\begin{solution}{iq:iv2:markets-and-microstructure:24}
La moyenne de la journée vaut 100.09. Le rapport affirme qu'on a payé 21~cents au-dessus. Le coût réel est d'environ 30~cents. Si l'on était à soi seul la totalité du volume, l'indice de référence égalerait le prix obtenu, et le rapport indiquerait zéro.

La moyenne vaut $0.7 \times 100 + 0.3 \times 100.30 = 100.09$. L'écart est $100.30 - 100.09 = 0.21$, soit le véritable 0.30 multiplié par $(1 - 0.3)$. L'indice de référence contient donc l'ordre lui-même, et il embellit le résultat d'un gros ordre.

\iqwrong{Utiliser cet indice de référence pour un gros ordre. Comparer aussi avec le prix d'arrivée\index{prix d'arrivée}, puisque celui-ci ne bouge pas quand on agit.}
\iqwhy{Un desk d'exécution rémunéré sur cet indice a intérêt à peser davantage dans le volume, ce qui peut aller à l'encontre de ce que veut le client.}
\end{solution}

\begin{solution}{iq:iv2:markets-and-microstructure:25}
On est indifférent pour $\lambda \approx 13.3$. Au-delà de ce seuil, on prend le prix du dealer.

Le dealer propose un montant certain de 0.8~\%. Travailler l'ordre coûte $0.5~\% + \lambda \times (1.5~\%)^2$ avec une pénalité. Donc $0.003 = \lambda \times 0.000225$, d'où $\lambda = 13.33$. Pour fixer les idées, les 0.3~\% économisés en travaillant l'ordre s'accompagnent de 1.5~\% de risque~: 0.2 écart-type de rendement moyen par unité de risque. Sur un ou deux jours, c'est un ratio de Sharpe élevé à encaisser sans aucune opinion de marché.

Il faut aussi se demander ce que révélera l'exécution chez le dealer~: le marché est-il au courant du bloc~? Détient-on d'autres titres du même type~? Votre estimation de 0.5~\% tient-elle compte de la fuite d'information~? Quelle est la corrélation de la position avec le reste du portefeuille~?

\iqwrong{Ne comparer que les coûts espérés. À eux seuls, travailler l'ordre l'emporte toujours, et personne ne paierait jamais pour un prix de risque.}
\end{solution}

\begin{solution}{iq:iv2:markets-and-microstructure:26}
On gagne avec une probabilité de 28.125~\%, et l'espérance de gain est d'environ 0.48~cent par action et par tentative.

On prend le carré unité, en unités de 200 microsecondes, avec votre temps $Y = 200 + 200U$ et celui de l'annulation $X = 150 + 200W$. On gagne si $Y < X$, donc si $U < W - \tfrac14$. Il s'agit d'un triangle de côté $\tfrac34$, dont l'aire vaut $\tfrac12 \times \tfrac{9}{16} = \tfrac{9}{32} = 0.28125$.

Quand on gagne, on touche 2~cents moins les 0.3~cent de frais, soit 1.7~cents. Donc $0.28125 \times 1.7 \approx 0.478$.

Une course de ce type n'a pas de gagnant net, puisque le côté le plus lent gagne plus d'un quart du temps. Réduire la variance de sa propre latence compte autant que sa moyenne.

\iqwrong{Dire que le côté le plus rapide gagne toujours. Les latences se dispersent, et une médiane plus lente gagne encore parfois.}
\end{solution}

\begin{solution}{iq:iv2:markets-and-microstructure:27}
Cela rapporte au-delà d'une probabilité d'exécution d'un tiers. À 40~\%, on gagne 0.1~cent par action.

Le gain vaut $p \times 1 - (1 - p) \times 0.5$, qui s'annule pour $p = \tfrac13$. Pour $p = 0.4$~: $0.4 - 0.3 = 0.1$.

Ce qui manque, c'est l'identité des contreparties du dark pool. Les exécutions dans un dark pool\index{dark pool} viennent souvent de ceux qui sont prêts à vous vendre juste avant une baisse, de sorte qu'une exécution au prix milieu n'est pas toujours une économie. Manque aussi le délai, ainsi que la question de savoir si la fuite d'information dépend du dark pool.

\iqwrong{Traiter le prix milieu comme une économie gratuite. Une transaction au prix milieu dans le dark pool peut être pire qu'une transaction au cours vendeur, si le prix baisse ensuite.}
\end{solution}

\begin{solution}{iq:iv2:markets-and-microstructure:28}
Le teneur de marché gagne 0.5~cent par action. Il peut payer pour le flux, car les ordres de particuliers ne sont en moyenne pas informés des mouvements à court terme, de sorte que la sélection adverse est faible.

La demi-fourchette vaut 1~cent. Il cède 0.2 d'amélioration, paie 0.1 et perd 0.2 au titre de la sélection adverse, soit $1 - 0.2 - 0.1 - 0.2 = 0.5$.

La cotation de la bourse doit couvrir tout le flux, y compris le flux informé. Un dealer qui peut choisir le flux avec lequel il négocie peut offrir un meilleur prix à ce flux. Cela soulève la question de savoir qui supporte la sélection adverse restée sur la bourse, et si le courtier doit envoyer l'ordre là où le prix du client est le meilleur.

\iqwrong{Dire que cela «~doit être du front-running~». L'économie tient à la seule faible teneur en information du flux.}
\end{solution}

\begin{solution}{iq:iv2:markets-and-microstructure:29}
Les rendements à la minute donnent environ 2.08~\%, soit 4~\% de trop. Les rendements à la seconde donnent environ 4.77~\%, soit plus du double de la vraie valeur. Il faudrait échantillonner à environ cinq minutes, ou utiliser un estimateur robuste au bruit, tel que la volatilité réalisée à deux échelles.

Chaque rendement échantillonné ajoute une variance de bruit $s^2/2$, où $s$ est l'écart exprimé en termes de rendement~: $(0.0004)^2/2 = 8 \times 10^{-8}$. Sommé sur 390 minutes, cela fait 0.0000312, auquel s'ajoute $0.02^2 = 0.0004$. La variance vaut donc 0.000431, et sa racine 2.08~\%. Sur 23\,400 secondes, le bruit vaut 0.00187, et le total 0.00227, soit une racine de 4.77~\%.

La vraie variance ne dépend pas de la fréquence d'échantillonnage, et le bruit croît en proportion directe avec elle. Le signature plot de la volatilité réalisée en fonction de l'intervalle d'échantillonnage\index{volatilité réalisée} le montre.

\iqwrong{Supposer que davantage de données est toujours préférable. Un échantillonnage plus fréquent donne ici une estimation plus précise du mauvais nombre.}
\end{solution}

\begin{solution}{iq:iv2:markets-and-microstructure:30}
Vendre pour 15~millions de dollars de l'indice. La volatilité résiduelle vaut environ 19.8~\%, soit environ 1.98~millions de dollars par an, et environ 124\,000~\$ par jour.

Le bêta vaut $\rho\sigma_s/\sigma_m = 0.75 \times 30/15 = 1.5$. Ce qui reste est $\sigma_s\sqrt{1 - \rho^2} = 30~\% \times \sqrt{0.4375} \approx 19.84~\%$. Par jour, cela fait $1.98/16$.

La couverture supprime 56~\% de la variance, mais seulement un tiers de la volatilité. Il faut mesurer le risque spécifique à l'action qui subsiste.

\iqwrong{Couvrir dollar pour dollar, avec 10~millions de dollars. Il reste alors un bêta de 0.5 au marché, et donc un risque de marché non couvert.}
\end{solution}

\begin{solution}{iq:iv2:markets-and-microstructure:31}
Le meilleur $\delta$ vaut 2~cents, avec un taux d'exécution de $e^{-1} \approx 0.37$ par seconde et un profit d'environ 0.74~cent par seconde. Doubler $k$ ramène $\delta$ à 1~cent et divise par deux le taux de profit, tandis que le taux d'exécution reste à 0.37.

Le taux de profit est $\delta e^{-k\delta}$, et sa dérivée s'annule en $\delta = 1/k$. Là, le taux de profit vaut $1/(ke)$. Ainsi, la concurrence, qui se traduit par un $k$ plus élevé, resserre la cotation et divise le profit par deux, les exécutions restant inchangées.

\iqwrong{Coter aussi serré que possible pour obtenir davantage d'exécutions. Chaque exécution rapporte moins, et l'optimum est un compromis entre les deux.}
\end{solution}

\begin{solution}{iq:iv2:markets-and-microstructure:32}
On s'attend à payer $-p + (1 - p) \times 1.8$ cents par rapport au prix milieu du jour, soit $1.8 - 2.8p$. Afficher un ordre à cours limité est préférable à croiser la fourchette si $p > 0.8/2.8 \approx 28.6~\%$. Il ne faut rien retrancher pour la baisse du prix milieu lors d'une exécution~: on a acheté au cours acheteur de toute façon.

Si l'ordre est exécuté, on paie 1~cent sous le prix milieu du jour. Sinon, on paie $0.8 + 1 = 1.8$ cents au-dessus. Il faut le comparer au croisement à $+1$. En général, $p^* = d/(2 + d)$, et pour $d = 0$, afficher un ordre est toujours préférable.

On détient l'action dans chaque branche, donc sa valeur future ne diffère pas entre les choix. La sélection adverse apparaît dans la taille de $d$~: une exécution se produit plus souvent quand le prix baisse, et une non-exécution coûte davantage quand il monte. Retrancher davantage reviendrait à la compter deux fois.

\iqwrong{Compter la sélection adverse à part sur la branche exécutée. Un teneur de marché, qui doit alors vendre, doit le faire. Un acheteur qui veut détenir l'action ne doit pas le faire.}
\end{solution}

\begin{solution}{iq:iv2:markets-and-microstructure:33}
Le VPIN\index{VPIN} vaut 0.35. Il est censé mesurer la part du flux unilatéral, comme approximation du flux toxique.

Les déséquilibres valent 20, 40, 0 et 80, soit une moyenne de 35 par tranche de 100.

Les mises en garde sont les suivantes. La classification des transactions est bruitée. Le déséquilibre peut refléter un gros ordre non informé, comme celui d'un fonds indiciel. Les valeurs élevées vont aussi de pair avec la volatilité~; il faut donc vérifier si le VPIN prédit quelque chose au-delà de la volatilité et du volume avant de dimensionner une décision qui s'appuie dessus.

\iqwrong{Lire une valeur élevée comme une preuve de trading informé. Il mesure un déséquilibre, et le reste relève de l'interprétation.}
\end{solution}

\begin{solution}{iq:iv2:markets-and-microstructure:34}
Un seul balayage coûte 55~\$. Les tranches coûtent 19~\$, et le risque est un écart-type d'environ 84~\$. Les tranches font donc économiser 36~\$ en coût espéré, tout en portant un risque d'environ 2.3 fois ce montant.

Un balayage coûte $1, 2, \ldots, 10$ cents sur chaque lot de 100 actions~: $100 \times 55$ cents. Les tranches coûtent $1 + 0.2(k - 1)$ cents sur la tranche $k$, soit une moyenne de 1.9~cents sur 1\,000 actions.

Le risque vient des actions que vous n'avez pas encore achetées~: 900 après la première minute, puis 800, et ainsi de suite jusqu'à 100. La somme des carrés vaut $100^2 \times 285$, donc la racine est $100 \times 16.9$, et à 5 cents la minute cela fait environ 84~\$.

\iqwrong{Ignorer le risque parce qu'il est de moyenne nulle. Une vision moyenne-variance le pondère, et si vous avez une vue à court terme, le coût est réel.}
\end{solution}

\begin{solution}{iq:iv2:markets-and-microstructure:35}
Vendre 625 d'abord et 375 ensuite, pour un coût total d'environ 143.75~\$. Un partage égal coûte 150~\$, et tout en une fois coûte 200~\$. Pour $c = 0$, le meilleur partage est l'égal partage. Pour $c = 0$ et $b = 2a$, tous les partages coûtent la même chose.

Posons $y = x_2$ et $x_1 = X - y$. Le coût vaut $a(X - y)^2 + ay^2 + b(X - y)y + cy^2$. Sa dérivée en $y$ est $-2a(X - y) + 2ay + b(X - 2y) + 2cy$, qui s'annule en
\[
y = \frac{X(2a - b)}{4a - 2b + 2c} = \frac{1000 \times 0.0003}{0.0008} = 375.
\]
Le coût vaut alors $0.0002(625^2 + 375^2) + 0.0001 \times 625 \times 375 + 0.0001 \times 375^2 = 106.25 + 23.44 + 14.06$.

Avec $c = 0$, $y = X/2$~: partager également répartit l'impact convexe, et le risque avance le partage. Avec $c = 0$ et $b = 2a$, le coût vaut $a(x_1 + x_2)^2$ pour tout partage, donc le calendrier ne sert à rien, et l'effet de prix de la première tranche annule le gain du partage.

\iqwrong{Penser que l'aversion au risque pousse à trader plus tard. Détenir est le risque, donc cela avance la vente.}
\end{solution}

\begin{solution}{iq:iv2:markets-and-microstructure:36}
La probabilité que la contrepartie vendeuse passe en premier est d'environ 91~\%, donc le mouvement espéré vaut environ $+0.41$ tick. Le microprix place la juste valeur 0.2 tick au-dessus du prix milieu. Les deux diffèrent parce que le microprix est une pondération linéaire simple, alors que la course est une probabilité de queue.

La contrepartie vendeuse passe en premier si 3 achats arrivent avant 7 ventes. Regardons les 9 premiers ordres, qui suffisent à trancher~: la contrepartie vendeuse passe en premier si au moins 3 sont des achats. Alors $P = 1 - (1 + 9 + 36)/512 = 466/512 \approx 91.0~\%$. Le mouvement espéré vaut $\tfrac12(0.910 - 0.090) = 0.41$ tick.

Les pondérations du microprix\index{microprix} sont $7/10$ et $3/10$, donc il se situe $0.7 - 0.5 = 0.2$ tick au-dessus du prix milieu. Par vérification, avec des files égales, les deux donnent 0.

La course surestime ce que l'on observe en pratique, car les files sont réalimentées. La réalité se situe en général entre les deux. Il faut donc ajuster le mouvement en fonction du déséquilibre à partir des données, et l'on constate en général qu'il est courbe.

\iqwrong{Donner les 70~\% tirés des tailles comme probabilité. C'est la part de taille, et non la chance de gagner une course jusqu'à l'épuisement de la file.}
\end{solution}

\begin{solution}{iq:iv2:markets-and-microstructure:37}
Un achat qui suit un achat déplace le prix de 0.2 cent, et une vente qui suit un achat le déplace de $-1.8$ cent. Après une série de cinq achats partant d'un état neutre, le prix a bougé de 1.8 cent. Le flux d'ordres est prévisible, donc le prix intègre déjà ce qui est attendu~: seule la surprise le fait bouger.

Le signe suivant attendu après un achat vaut 0.8. Un achat constitue une surprise de 0.2, et une vente une surprise de $-1.8$. Chacun de ces mouvements a une moyenne nulle, puisque $0.9 \times 0.2 + 0.1 \times (-1.8) = 0$, avec la probabilité d'achat $(1 + 0.8)/2 = 0.9$. Les variations de prix ne sont donc pas autocorrélées, bien que les signes le soient.

Sur une série de cinq, la première est une surprise de 1 et les quatre autres de 0.2 chacune, soit $1 + 0.8$. L'impact d'une suite d'achats est concave en longueur, et c'est la même concavité que l'on observe sur les gros ordres. Il s'agit d'une forme simple du modèle de Madhavan--Richardson--Roomans\index{Madhavan-Richardson-Roomans model}.

\iqwrong{Affirmer que la pression acheteuse doit continuer à pousser le prix à la hausse parce que les ordres sont autocorrélés. Des prix aussi faciles à prévoir seraient arbitrés aussitôt.}
\iqfollow{Et si c'est vous qui découpez l'ordre~?}{Le marché lit votre série comme une continuation probable, et il vous facture donc ce risque dans le prix. Cela rend vos premières tranches chères et les suivantes bon marché.}
\end{solution}

\begin{solution}{iq:iv2:markets-and-microstructure:38}
Il faut 2.5 jours. Votre coût d'impact est d'environ 0.94~\%. L'écart-type est d'environ 1.83~\%. Vous cotez donc une décote d'environ 1.86~\%. Une couverture par un indice ramène l'écart-type à environ 1.10~\%, et la moitié de cette valeur porte la décote à environ 1.49~\%.

Le rythme vaut $0.2 \times 2$ millions par jour, donc 1 million prend 2.5 jours. L'impact vaut $\tfrac23 \times 2~\% \times \sqrt{0.5} \approx 0.943~\%$.

Pour le risque, vous détenez une action dont la position décroît linéairement de 1 à 0 sur $T$ jours. Votre variance vaut $\sigma^2\int_0^T (1 - t/T)^2 dt = \sigma^2T/3$. L'écart-type vaut donc $2~\% \times \sqrt{2.5/3} \approx 1.83~\%$. Avec une couverture, $1.2~\% \times \sqrt{2.5/3} \approx 1.10~\%$.

Le total vaut $0.943 + \tfrac12 \times 1.826 \approx 1.86~\%$. Avec la couverture, il vaut $0.943 + \tfrac12 \times 1.095 \approx 1.49~\%$. Il faut ensuite ajouter le risque que ce client sache quelque chose~: s'il est informé, vous perdez quel que soit le contenu de l'information.

\iqwrong{Ne facturer que l'impact. Le teneur de marché n'a pas de vue sur l'action, et pour les 2.5 jours de risque, il faut être rémunéré.}
\end{solution}

\begin{solution}{iq:iv2:markets-and-microstructure:39}
Le prix d'indifférence côté achat vaut $s - \gamma\sigma^2\tau(q + \tfrac12)$, et celui côté vente $s - \gamma\sigma^2\tau(q - \tfrac12)$. Leur moyenne, le prix de réservation\index{prix de réservation}, vaut $s - q\gamma\sigma^2\tau$. Avec ces chiffres, il vaut 99, et l'on cote donc environ 98.25 et 99.75.

Avec une utilité exponentielle et un prix gaussien, l'équivalent certain de la détention de $q$ actions vaut $qs - \tfrac12\gamma q^2\sigma^2\tau$. Acheter une action de plus à $p$ vaut le coup si $(q + 1)s - \tfrac12\gamma(q + 1)^2\sigma^2\tau - p \geq qs - \tfrac12\gamma q^2\sigma^2\tau$. On obtient $p \leq s - \gamma\sigma^2\tau(q + \tfrac12)$. Vendre une action donne l'autre de la même façon.

Avec les chiffres, $\gamma\sigma^2\tau = 0.1 \times 4 \times 0.5 = 0.2$, et le prix de réservation vaut $100 - 5 \times 0.2 = 99$. Le spread vaut $0.2 + 20\ln(1 + 0.1/1.5) \approx 1.49$, donc sa moitié vaut 0.745. Ce sont les cotations d'Avellaneda--Stoikov\index{Avellaneda-Stoikov model}.

Si $q = 0$, les cotations sont centrées sur $s$. Si vous êtes long, vous les abaissez pour vendre et pour acheter moins.

\iqwrong{N'incliner qu'un seul côté. Les deux cotations se déplacent du même montant, ce qui conserve le même spread avec un nouveau centre.}
\end{solution}

\begin{solution}{iq:iv2:markets-and-microstructure:40}
Non. Le premier cas est du spoofing\index{spoofing}~: placer une offre dans l'intention de l'annuler avant exécution. Le second peut relever d'une stratégie licite dont les ordres changent effectivement. Le droit américain interdit le premier par son intention. Le Commodity Exchange Act, à 7 U.S.C. 6c(a)(5)(C), interdit de placer une offre ou une demande dans l'intention de l'annuler avant son exécution.

Un taux d'annulation ne peut donc pas trancher. Dans l'affaire contre un trader, une cour d'appel américaine a rejeté l'argument selon lequel la règle serait floue parce que les traders à haute fréquence annulent 98~\% de leurs ordres. La question porte sur la finalité de l'ordre. Les grosses offres du premier algorithme ne sont placées que pour déplacer les prix des autres, donc elles ne visent pas à trader, alors que celles du second visaient à trader au prix alors affiché.

Pour le détecter, regardez les ordres où une grosse taille au repos d'un côté coïncide avec des exécutions de l'autre côté, et où les gros ordres sont annulés dès que ces exécutions surviennent. Comparez la distance des ordres au marché, le ratio ordres sur transactions par configuration, et la part des ordres du côté opposé qui s'exécutent. Transmettez au service conformité dès que vous identifiez une configuration. Ne jugez pas seul de l'intention.

\iqwrong{«~Nous annulons tout, donc nous sommes tous des spoofers~» ou «~nous n'avions pas l'intention de nuire, donc c'est correct~». La loi recherche une intention d'annuler au moment de l'ordre, et elle n'exige pas de préjudice.}
\iqfollow{Cela ne concerne-t-il que les contrats à terme~?}{La disposition citée vise les marchés relevant de cette loi, et les marchés de titres et les autres ont leurs propres règles de manipulation. Il faut savoir quel régime couvre chaque desk.}
\end{solution}

\begin{solution}{iq:iv2:markets-and-microstructure:41}
Ma réponse recommandée~: les investisseurs de long terme aux petits ordres gagnent probablement à un spread plus étroit, ceux qui passent de gros ordres peuvent ne pas y gagner, les teneurs de marché en tête de file perdent, ceux en queue peuvent gagner, et les traders rapides qui peuvent passer devant gagnent. L'effet net dépend de l'écart entre le spread naturel et le tick.

Le tick constitue un plancher pour le spread. Si le spread naturel est inférieur à 1 cent, diviser le tick par deux le resserre. Un petit ordre paie alors moins.

Mais un tick protège aussi une file. Aujourd'hui, pour passer devant une longue file, il faut améliorer d'un cent entier, ce qui coûte cher. Avec un tick d'un demi-cent, passer devant coûte deux fois moins. La grosse taille affichée qui se trouve là aujourd'hui vaut donc moins, et la profondeur baisse au meilleur cours. Un gros ordre peut désormais devoir payer davantage de niveaux, il peut donc gagner moins, voire perdre.

Ceux qui sont en tête de file perdent une position pour laquelle ils ont payé. Ceux qui sont en queue trouvent la file plus courte. Les traders rapides peuvent reprix à bas coût, et ce sont eux qui gagnent le plus à passer devant.

Pour en juger, mesurer les spreads cotés et effectifs, la profondeur au meilleur cours et à proximité, le coût des gros ordres au travers de l'implementation shortfall, les spreads réalisés et la longueur des files, ainsi que la sélection adverse subie par le teneur de marché. Comparer avec des actions semblables non modifiées, avant et après.

\iqwrong{«~Un tick plus petit signifie un spread plus étroit, donc tout le monde y gagne.~» Le spread n'est qu'un coût, et la profondeur comme la valeur des files changent aussi.}
\end{solution}

\begin{solution}{iq:iv2:markets-and-microstructure:42}
La demi-fourchette de seuil de rentabilité vaut 0.5 cent. Avec une seule entreprise, elle vaut 0. Avec un très grand nombre, elle tend vers environ 0.67 cent. Supprimer la course supprimerait cette composante du spread.

Les transactions ordinaires paient $h$ à un rythme de 1 par seconde. Un signal arrive à un rythme de 0.5. Avec trois entreprises, le teneur de cotation est le premier à annuler une fois sur trois, et il se fait donc prendre au piège deux fois sur trois. Chaque fois, il perd $2 - h$. Donc
\[
1 \times h = 0.5 \times \tfrac23 \times (2 - h),
\]
ce qui donne $h = \tfrac23/\tfrac43 = 0.5$. Avec $N$ entreprises, le facteur vaut $(N - 1)/N$~: pour $N = 1$ le terme s'annule, et à mesure que $N$ croît il tend vers 1, donc $h \to 1/1.5$.

C'est le résultat du travail de Budish, Cramton et Shim sur la course à la vitesse\index{latency arbitrage}~: la course elle-même est le coût, et ce sont les traders ordinaires qui le paient dans le spread. Les enchères par lots fréquentes ont été proposées pour le supprimer. Être plus rapide ne le réduit pas~: si tout le monde est plus rapide, l'égalité à trois demeure.

\iqwrong{Affirmer que l'investissement dans la latence resserre les spreads parce que les prix se mettent à jour plus vite. Ici $h$ ne dépend pas de la rapidité des entreprises, seulement de leur nombre.}
\end{solution}

\begin{solution}{iq:iv2:markets-and-microstructure:43}
Vous espérez gagner $50q/(1000 + q)$ par transaction. Ce gain croît avec $q$ et se stabilise vers 50. Si vous pouvez prendre au plus 200 lots en une transaction, montrez qu'il faut environ 333 lots.

Votre part de toute transaction vaut $q/(1000 + q)$. Sur une petite transaction, vous gagnez un demi-tick sur 200 lots. Sur une grande, vous gagnez un demi-tick et perdez un tick sur 800 lots. Donc le montant par unité de part vaut $0.9 \times 200 \times 0.5 + 0.1 \times 800 \times (0.5 - 1) = 90 - 40 = 50$.

Ainsi, à $q = 100$ on gagne environ 4.5, à $q = 1\,000$ on gagne 25, et à $q = 3\,000$ on gagne 37.5. Une taille plus grande aide toujours, en espérance, donc tout le monde surdimensionne. Ce qui limite, c'est le risque~: la plus grosse exécution vaut $800q/(1000 + q)$, et l'égaler à 200 donne $q = 1000/3$.

Les autres voient la même chose, donc leur taille augmente aussi, et le 1\,000 n'est pas fixe. C'est pourquoi les carnets en pro-rata affichent bien plus que ce que les gens veulent.

\iqwrong{Afficher ce que l'on veut négocier. En pro-rata, on n'est exécuté que pour une partie de son affichage, donc on affiche davantage, et l'on risque d'être exécuté pour la totalité.}
\end{solution}

\begin{solution}{iq:iv2:markets-and-microstructure:44}
Ce n'est pas significatif~: l'erreur type de la différence vaut environ 2.12~points de base, donc $t \approx 0.57$. Il faudrait environ 9\,800 ordres par courtier. Le coût est si bruité par rapport à l'écart qu'on ne peut pas classer les courtiers sur quelques centaines d'ordres.

L'erreur type vaut $30 \times \sqrt{2/400}$. Pour l'échantillon requis, $n = 2(1.96 + 0.84)^2(30/1.2)^2 = 2 \times 7.84 \times 625 = 9\,800$.

Deux moyens de faire mieux. Utiliser un contrôle~: comparer le coût de chaque ordre à un modèle de ce qu'il devrait coûter, compte tenu de la taille, de la volatilité et de la fourchette, permet d'éliminer une grande part de la variance. Et équilibrer les ordres selon leur difficulté, pour que la chance dans le mélange ne tranche pas. L'affectation aléatoire de l'énoncé est la partie précieuse, car sans elle, un courtier peut recevoir les ordres difficiles.

\iqwrong{Choisir le courtier le moins cher sur 400 ordres, puis être surpris le trimestre suivant. Sur un échantillon de cette taille, le classement est proche d'un pile ou face.}
\end{solution}

\begin{solution}{iq:iv2:markets-and-microstructure:45}
La meilleure valeur de $y$ pour le prédateur est 25\,000 actions, soit la moitié d'une tranche, pour un profit de 625~\$. Votre coût supplémentaire est de 2\,500~\$, soit un tiers de votre coût de 7\,500~\$ sans prédateur. Un calendrier prévisible coûte donc bien plus au vendeur qu'il ne rapporte au prédateur, et il faut randomiser son timing et sa taille et dissimuler son volume total.

Sans prédateur, on vend $x = 50\,000$ à $p_0 - bx$, puis $x$ à $p_0 - 2bx$. Le coût vaut donc $3bx^2 = 3 \times 2\,500 = 7\,500$, puisque $bx^2 = 2\,500$.

Le prédateur vend $y$ à $p_0 - b(x + y)$, et rachète à $p_0 - 2bx$. Le profit vaut $y \cdot b(x - y)$, et il est maximal en $y = x/2$, où il vaut $bx^2/4 = 625$.

On reçoit moins de $by$ sur chacune de ses tranches, car la quantité $y$ du prédateur reste dans le prix jusqu'au rachat. Cela fait $2bxy = bx^2 = 2\,500$. Le reste de la perte revient à ceux qui ont acheté auprès de soi et du prédateur, à des prix plus bas.

\iqwrong{Penser qu'un ordre étalé dans le temps n'a aucune fuite parce que chaque tranche est petite. Un schéma prévisible est lui-même la fuite.}
\end{solution}

\begin{solution}{iq:iv2:markets-and-microstructure:46}
L'horizon optimal vaut environ 0.58 jour. Là, le coût et le risque valent chacun environ 0.88~\%, soit un total d'environ 1.75~\%. La loi empirique en racine carrée dit que le coût dépend de $Q/V$ avec peu de dépendance au temps, donc le gain que ce modèle attribue à une exécution lente peut être surestimé.

On écrit l'objectif sous la forme $aT^{-1/2} + bT^{1/2}$, avec $a = \tfrac23 \times 2~\% \times \sqrt{0.25} \approx 0.667~\%$ et $b = 2~\%/\sqrt 3 \approx 1.155~\%$. La dérivée donne $T = a/b \approx 0.577$. En ce point, les deux termes sont égaux, chacun valant $\sqrt{ab} \approx 0.877~\%$.

Là où la loi est vraie, le coût baisse à peine quand l'horizon s'allonge, tant que les ordres restent dans son domaine. Aller lentement n'apporte alors que du risque, et la bonne réponse est plus courte. Il faut vérifier laquelle s'applique à ses propres données~: on regroupe les ordres par durée et l'on regarde si le coût à $Q/V$ donné en dépend.

\iqwrong{Dire qu'une exécution plus lente est toujours moins chère. Un étalement plus long ajoute du risque de prix et de la fuite. Et les données ne montrent pas de gains importants à étirer un ordre fixe.}
\end{solution}

\begin{solution}{iq:iv2:markets-and-microstructure:47}
Les bandes sont à 95~\$ et 105~\$. Si le marché reste au-delà de la bande pendant 15 secondes en état de limite, la bourse de cotation déclare une pause d'au moins 5 minutes. Les bandes sont doublées pendant les 15 premières minutes de la séance, et pendant les 25 minutes qui précèdent la clôture. Pendant la pause, on ne peut pas couvrir son delta sur cette action, et l'on reste exposé. On se couvre donc avec un instrument lié qui continue de se négocier, ou l'on réduit sa taille.

Dans le plan Limit Up--Limit Down\index{limit up-limit down}, les bandes sont fixées autour du prix moyen des transactions des cinq dernières minutes. Une transaction ne peut pas être exécutée en dehors de la bande. En état de limite, le cours vendeur est à la bande inférieure et il n'y a pas de cours acheteur à ce niveau, donc les vendeurs ne trouvent pas d'acheteurs à un prix autorisé, et le prix ne peut pas baisser davantage pour le moment. Au bout de 15 secondes, la négociation est suspendue. La reprise se fait par enchère.

On a vendu un put, et on s'est couvert en vendant l'action à découvert. Quand l'action baisse, le delta du put augmente et il faut vendre davantage. Pendant la pause, on ne peut pas le faire, et la reprise peut se faire avec un gap nettement sous la bande, de sorte que l'on subit une perte de gamma que l'on n'a pas pu couvrir.

On utilise des contrats à terme sur indice ou un fonds sectoriel comme couverture partielle, et l'on accepte la base. Il faut vérifier comment la plateforme où l'on négocie l'option traite la pause avant de s'y fier. Un desk raisonnable limite aussi le short gamma sur les valeurs les plus susceptibles d'être suspendues, comme celles qui publient des résultats.

\iqwrong{Penser que les bandes protègent d'un gap. Elles arrêtent les transactions à un prix. Elles n'empêchent pas le cours de s'ouvrir plus bas.}
\end{solution}

\begin{solution}{iq:iv2:markets-and-microstructure:48}
L'initié négocie $x = \beta(v - p_0)$ avec $\beta = \sigma_u/\sigma_v$, et le prix réagit avec $\lambda = \sigma_v/(2\sigma_u)$. Avec les valeurs numériques, $\beta = 500$ actions par dollar, $\lambda = 0.001$ dollar par action, le profit espéré de l'initié est de 1\,000~\$, et la variance restante vaut 2. Ce sont les noise traders qui paient.

Sachant $v$, l'initié choisit $x$ pour maximiser $E[x(v - p_0 - \lambda(x + u))] = x(v - p_0) - \lambda x^2$. Donc $x = (v - p_0)/(2\lambda)$ et $\beta = 1/(2\lambda)$.

Le teneur de marché fixe son prix par la régression de $v$ sur le flux d'ordres $y = x + u$~: $\lambda = \mathrm{Cov}(v, y)/\mathrm{Var}(y) = \beta\sigma_v^2/(\beta^2\sigma_v^2 + \sigma_u^2)$. En remplaçant $\beta = 1/(2\lambda)$, il vient $\beta^2\sigma_v^2 = \sigma_u^2$. Donc l'ordre de l'initié a la même variance que le bruit, et $\lambda = \sigma_v/(2\sigma_u)$.

Le profit espéré vaut $E[(v - p_0)^2]/(4\lambda) = \sigma_v\sigma_u/2 = 2 \times 1000/2 = 1\,000$. La variance de $v$ sachant $y$ vaut $\sigma_v^2(1 - \tfrac12) = 2$~: la moitié de l'information entre dans le prix.

Dans ce modèle de Kyle\index{modèle de Kyle}, le teneur de marché ne gagne ni ne perd, en moyenne. Le profit de l'initié est donc égal à la perte des noise traders. Doubler le bruit divise $\lambda$ par deux et double le profit de l'initié, tout en laissant le prix aussi informatif.

\iqwrong{Penser qu'un bruit accru rend le marché moins informatif. L'initié négocie simplement davantage, de sorte que la moitié exacte de la variance passe toujours dans le prix.}
\iqwhy{C'est d'où vient $\lambda$, la variation de prix par action, comme mesure de la liquidité, et pourquoi on le lit comme une affirmation sur l'information.}
\end{solution}

\begin{solution}{iq:iv2:markets-and-microstructure:49}
La trajectoire optimale est $x(t) = X\sinh(\kappa(T - t))/\sinh(\kappa T)$. À mi-parcours, on a vendu environ 55.7~\% pour $\kappa T = 1$, et environ 78.7~\% pour $\kappa T = 3$. Par rapport à une vente à rythme constant, le coût augmente d'environ 1.9~\% et la variance baisse d'environ 11.7~\% pour $\kappa T = 1$. Pour $\kappa T = 3$, le coût augmente d'environ 55~\%, et la variance baisse d'environ 51~\%. Pour $\lambda = 0$, la trajectoire est une droite.

Avec $v = -\dot x$, l'intégrande de l'objectif vaut $L = \eta\dot x^2 + \lambda\sigma^2x^2$. L'équation d'Euler-Lagrange\index{calcul des variations} $2\eta\ddot x = 2\lambda\sigma^2 x$ donne $\ddot x = \kappa^2x$. Avec $x(0) = X$ et $x(T) = 0$, on obtient la forme en $\sinh$ ci-dessus. Lorsque $\kappa \to 0$, elle tend vers $X(1 - t/T)$.

La part vendue à mi-parcours est $1 - \sinh(\kappa T/2)/\sinh(\kappa T)$~: $1 - 0.5211/1.1752$ et $1 - 2.1293/10.018$.

Avec $T = 1$, l'intégration donne
\[
\int_0^1\dot x^2\,dt = X^2\kappa^2\frac{\tfrac12 + \sinh(2\kappa)/(4\kappa)}{\sinh^2\kappa},\qquad
\int_0^1 x^2\,dt = X^2\frac{\sinh(2\kappa)/(4\kappa) - \tfrac12}{\sinh^2\kappa}.
\]
Pour une vente à rythme constant, ces intégrales valent $X^2$ et $X^2/3$. Ainsi, pour $\kappa = 1$~: 1.019 et 0.2945. Pour $\kappa = 3$~: 1.552 et 0.1625.

C'est le modèle d'Almgren--Chriss\index{modèle d'Almgren--Chriss}. La trajectoire est raide au départ, car détenir la position est le risque. Le compromis se voit le mieux sur la frontière efficiente~: un petit écart par rapport au cas linéaire procure une forte réduction du risque pour un faible coût, et les écarts ultérieurs sont coûteux.

\iqwrong{Vendre uniformément et appeler cela l'optimum. C'est optimal seulement sans aversion au risque. Tout $\lambda > 0$ avance les ventes.}
\iqfollow{Que vous apprend $\kappa$~?}{Son inverse est le temps caractéristique de l'ordre~: la durée approximative qu'un ordre prend pour un trader avec ce coût et ce risque. Si $T$ est bien supérieur, l'essentiel de l'ordre est exécuté dès le début.}
\end{solution}

\begin{solution}{iq:iv2:markets-and-microstructure:50}
Les tailles sont proportionnelles à $1$, $1 - a$ et $1$. Avec $a = 0.5$, les parts valent 40~\%, 20~\% et 40~\%, pour un coût de $0.3\kappa X^2$. Cela fait environ 1.8~\% de moins que trois tranches égales (coût d'environ $0.3056\kappa X^2$) et 40~\% de moins que tout en une fois ($0.5\kappa X^2$). Quand $a \to 0$, les tailles tendent vers des tiers égaux. Quand $a \to 1$, la transaction du milieu disparaît, et l'on échange la moitié à chaque extrémité.

On écrit le coût sous la forme $\tfrac12\kappa\,x^\top Gx$, avec $G$ valant 1 sur la diagonale, $a$ à côté et $a^2$ aux coins. L'optimum sous $\sum_i x_i = X$ a $Gx$ constant. Essayons $w = (1, 1 - a, 1)$. La première ligne donne $1 + a(1 - a) + a^2 = 1 + a$. La ligne du milieu donne $a + (1 - a) + a = 1 + a$. Donc $Gw = (1 + a)\mathbf 1$, et l'optimum est proportionnel à $w$.

Avec le total $S = 3 - a$, le coût vaut $\tfrac12\kappa X^2(1 + a)/S$. Pour $a = 0.5$~: $\tfrac12 \times 1.5/2.5 = 0.3$. Trois tranches égales coûtent $\tfrac12\kappa(X/3)^2(3 + 4a + 2a^2) = \tfrac12\kappa X^2 \times 5.5/9$.

C'est la forme du modèle d'Obizhaeva--Wang\index{modèle d'Obizhaeva--Wang}~: de grandes transactions au début et à la fin, et un écoulement plus lent entre les deux. La transaction du début dispose de tout le temps pour se résorber. La transaction de la fin ne laisse rien derrière elle qui puisse nuire, donc elle peut être importante.

\iqwrong{Supposer qu'un calendrier régulier doit être le meilleur. Avec un carnet qui se reconstitue, la première et la dernière transactions sont particulières, et le modèle dit d'en profiter.}
\iqfollow{Et si la décroissance n'est pas exponentielle~?}{Le résultat dépend de la forme de la décroissance. Avec une décroissance en loi de puissance, ou si l'impact est non linéaire, le calendrier change, et certaines formes permettent un aller-retour profitable, signe que le modèle n'est pas sûr.}
\end{solution}
